% === 第四章补充内容 ===

% --- 4.2 驱动函数分析、替换与编译模块实现 ---

\section{驱动函数分析、替换与编译模块实现}
\label{sec:builder_impl}

\subsection{LangGraph 智能体工作流实现}

\subsubsection{AgentState 数据结构}

系统使用 LangGraph 框架的 MessagesState 作为基础状态类，并扩展为包含特定功能的状态类型。AgentState 的定义如下：

\begin{lstlisting}[style=Python,abovecaptionskip=0pt,belowcaptionskip=0pt]
class AgentState(MessagesState):
    """智能体状态数据结构"""
    # 最终的结构化响应
    final_response: FunctionClassifyResponse
    # 当前处理的函数名
    function_name: str
\end{lstlisting}

该状态类继承自 MessagesState，自动包含 messages 字段用于维护对话历史。final\_response 字段存储 LLM 生成的结构化分类结果，function\_name 字段用于标识当前正在分析的函数。

\subsubsection{工作流构建与执行}

系统采用异步方式构建和执行 LangGraph 工作流。build\_graph() 函数负责创建工作流图，并配置所有必要的组件。

\begin{lstlisting}[style=Python,abovecaptionskip=0pt,belowcaptionskip=0pt]
async def build_graph():
    """构建函数分类智能体工作流"""
    global _graph, _client, _graph_db_path
    db_path = getattr(globs, "db_path", None) or ""
    
    # 缓存机制：如果数据库路径变化，重建工作流
    if _graph is not None and _graph_db_path != db_path:
        _graph = None
        _graph_db_path = None
    
    if _graph is not None:
        return _graph

    # 创建 MCP 客户端，连接到 CodeQL 服务器
    _client = MultiServerMCPClient({
        "lcmhal_collector": {
            "command": "python",
            "args": [
                "-m", "tools.collector.mcp_server",
                "--db-path", db_path,
                "--transport", "stdio"
            ],
            "transport": "stdio"
        },
    })
    
    # 获取工具列表（包括 CodeQL 查询工具和验证工具）
    from tools.builder.tool import verify_replacement as verify_replacement_tool
    from agents.builder_fixer_agent import builder_fixer_agent
    tools = await _client.get_tools()
    tools = tools + [verify_replacement_tool, builder_fixer_agent]
    
    # 绑定工具到模型
    model_with_tools = model.bind_tools(tools)
    # 设置结构化输出
    model_with_structured_output = model.with_structured_output(FunctionClassifyResponse)
    
    # 创建工具节点
    tool_node = ToolNode(tools)
    
    # 定义工作流节点函数
    async def agent_node(state: AgentState):
        """推理节点：LLM 分析并决策"""
        messages = state["messages"]
        response = await model_with_structured_output.ainvoke(messages)
        return {"messages": [response], "final_response": response}
    
    async def should_continue(state: AgentState) -> Literal["tools", "respond", "__end__"]:
        """判断是否需要继续调用工具"""
        messages = state["messages"]
        last_message = messages[-1]
        
        # 如果 LLM 请求调用工具
        if hasattr(last_message, "tool_calls") and last_message.tool_calls:
            return "tools"
        # 如果有最终响应
        elif "final_response" in state and state["final_response"]:
            return "respond"
        else:
            return "__end__"
    
    # 构建图
    workflow = StateGraph(AgentState)
    workflow.add_node("agent", agent_node)
    workflow.add_node("tools", tool_node)
    workflow.add_node("respond", lambda state: state)
    
    # 添加边和条件边
    workflow.add_edge(START, "agent")
    workflow.add_conditional_edges(
        "agent",
        should_continue,
        {
            "tools": "tools",
            "respond": "respond",
            "__end__": END,
        }
    )
    workflow.add_edge("tools", "agent")
    workflow.add_edge("respond", END)
    
    _graph = workflow.compile()
    _graph_db_path = db_path
    return _graph
\end{lstlisting}

工作流的执行流程如图~\ref{fig:workflow-execution}所示。系统从 START 状态开始，首先进入 agent 节点进行 LLM 推理。如果 LLM 决定调用工具，转移到 tools 节点执行工具调用，然后返回 agent 节点继续推理。如果 LLM 给出最终答案，转移到 respond 节点进行结果汇总，最后进入 END 状态。

\begin{figure}[htbp]
\centering
\begin{tikzpicture}[
    node distance=2cm,
    state/.style={rectangle, draw, rounded corners, minimum width=2.5cm, minimum height=1cm, align=center, font=\small},
    start/.style={ellipse, draw, fill=green!20, minimum width=1.5cm},
    end/.style={ellipse, draw, fill=red!20, minimum width=1.5cm},
    arrow/.style={->, thick, >=stealth}
]
    \node[start] (start) {START};
    \node[state, fill=blue!20] (agent) {Agent\\LLM推理};
    \node[state, fill=orange!20, right of=agent] (tools) {Tools\\工具执行};
    \node[state, fill=yellow!20, right of=agent, xshift=1cm] (respond) {Respond\\结果汇总};
    \node[end, right of=respond] (end) {END};
    
    \draw[arrow] (start) -- (agent);
    \draw[arrow] (agent) -- node[above] {有工具调用?} (decision);
    \draw[arrow] (decision) -- node[above] {是} (tools);
    \draw[arrow] (decision) -- node[above] {否} (respond);
    \draw[arrow] (tools) |- (agent);
    \draw[arrow] (respond) -- (end);
\end{tikzpicture}
\caption{LangGraph 工作流执行流程}
\label{fig:workflow-execution}
\end{figure}

\subsection{提示词工程}

\subsubsection{函数分类提示词}

函数分类提示词采用系统提示（system prompting）的方式，定义了完整的角色、任务、分类标准和输出格式。提示词的核心组件包括：

\textbf{（1）角色定义}

明确 LLM 的角色为嵌入式软件工程师，任务是替换驱动库函数以消除外设硬件依赖。

\textbf{（2）分类优先级}

定义了严格的分类优先级顺序：CORE > RECV > IRQ > INIT > LOOP > RETURNOK > SKIP > NEEDCHECK > NODRIVER。这个优先级确保了分类的一致性和可预测性。

\textbf{（3）可用工具列表}

详细列出了所有可用的工具及其使用场景：
\begin{itemize}
    \item \texttt{GetFunctionInfo}: 获取函数源码和基本信息
    \item \texttt{GetMMIOFunctionInfo}: 检查函数是否包含 MMIO 操作
    \item \texttt{GetStructOrEnumInfo}: 获取数据结构或枚举信息
    \item \texttt{GetFunctionCallStack}: 获取函数调用关系
    \item \texttt{GetDriverInfo}: 获取驱动模块信息
    \item \texttt{VerifyReplacement}: 验证替换代码
    \item \texttt{FixFunctionBuildErrors}: 委托编译错误修复给子智能体
\end{itemize}

\textbf{（4）代码生成规则}

引用外部的 \texttt{code\_generation\_rules.md} 文件，定义了严格的代码生成约束：

\begin{itemize}
    \item \textbf{信息源限制}：只能使用工具返回或项目头文件中存在的类型、成员、宏或函数
    \item \textbf{函数签名一致性}：必须与原始函数或头文件声明完全匹配
    \item \textbf{无伪造}：不能发明类型、结构体成员、宏、地址或全局变量
    \item \textbf{HAL/结构体成员}：只能使用原始函数或目标芯片 HAL 中存在的成员
    \item \textbf{其他不变量}：不改变函数返回类型，不修改 OS 相关调用
\end{itemize}

\subsubsection{Builder Agent 提示词}

Builder Agent 负责项目构建和错误修复，其提示词重点在于：

\textbf{（1）构建环境配置}：确保使用正确的编译工具链和构建脚本

\textbf{（2）错误解析能力}：能够理解 GCC/Clang 的错误输出格式

\textbf{（3）增量构建策略}：只重新编译修改过的文件，提高效率

\textbf{（4）集成验证流程}：构建成功后调用 VerifyReplacement 进行验证

\subsubsection{Fixer Agent 提示词}

Fixer Agent 负责修复编译错误，其提示词包括：

\textbf{（1）错误分类}：识别错误类型（语法错误、类型错误、链接错误）

\textbf{（2）修复策略}：针对不同错误类型采用不同的修复方法

\textbf{（3）保持一致性}：修复后的代码必须遵循代码生成规则

\textbf{（4）验证机制}：修复后重新调用 VerifyReplacement 确保问题解决

\subsection{替换策略的具体实现}

\subsubsection{Skip 策略实现}

Skip 策略适用于初始化类和控制类函数，实现逻辑如下：

\begin{lstlisting}[style=CStyle,abovecaptionskip=0pt,belowcaptionskip=0pt]
/**
 * Skip 策略：跳过硬件副作用
 * 适用场景：初始化函数、控制函数
 *
 * 实现要点：
 * 1. 移除所有 MMIO 寄存器访问
 * 2. 保留函数签名和返回值
 * 3. 保留状态变量更新
 * 4. 添加注释说明跳过的操作
 */

// 原始代码示例
HAL_StatusTypeDef HAL_UART_Enable(UART_HandleTypeDef *huart) {
    // 使能 UART 时钟
    __HAL_RCC_USART1_CLK_ENABLE();
    // 配置 UART 寄存器
    huart->Instance->CR1 |= USART_CR1_UE;
    // 设置状态
    huart->gState = HAL_UART_STATE_READY;
    return HAL_OK;
}

// Skip 策略替换代码
HAL_StatusTypeDef HAL_UART_Enable(UART_HandleTypeDef *huart) {
    // [SKIP] 硬件寄存器配置在仿真中不需要
    (void)huart;  // 避免未使用参数警告
    
    // 保留状态更新
    huart->gState = HAL_UART_STATE_READY;
    return HAL_OK;
}
\end{lstlisting}

\subsubsection{Strip 策略实现}

Strip 策略适用于发送类函数，需要保留状态管理逻辑但移除寄存器操作：

\begin{lstlisting}[style=CStyle,abovecaptionskip=0pt,belowcaptionskip=0pt]
/**
 * Strip 策略：剥离硬件代码、保留控制逻辑
 * 适用场景：发送数据函数
 *
 * 实现要点：
 * 1. 识别并删除所有寄存器读写操作
 * 2. 保留状态变量更新（gState, TxXferCount等）
 * 3. 保留错误处理逻辑
 * 4. 使用虚拟设备接口替代数据传输
 */

// 原始代码示例
HAL_StatusTypeDef HAL_UART_Transmit(UART_HandleTypeDef *huart,
                                     uint8_t *pData, uint16_t Size) {
    huart->gState = HAL_UART_STATE_BUSY_TX;
    for (int i = 0; i < Size; i++) {
        // [STRIP] 寄存器写入将被移除
        huart->Instance->DR = pData[i];
        // 轮询发送完成标志
        while (!(huart->Instance->SR & USART_SR_TXE));
    }
    huart->gState = HAL_UART_STATE_READY;
    return HAL_OK;
}

// Strip 策略替换代码
HAL_StatusTypeDef HAL_UART_Transmit(UART_HandleTypeDef *huart,
                                     uint8_t *pData, uint16_t Size) {
    huart->gState = HAL_UART_STATE_BUSY_TX;
    
    // [STRIP] 使用虚拟设备接口替代寄存器写入
    HAL_BE_Out(pData, Size);  // 将数据输出到虚拟设备
    
    huart->gState = HAL_UART_STATE_READY;
    return HAL_OK;
}
\end{lstlisting}

\subsubsection{Redirect 策略实现}

Redirect 策略适用于接收类函数，将操作重定向到仿真平台的虚拟设备接口：

\begin{lstlisting}[style=CStyle,abovecaptionskip=0pt,belowcaptionskip=0pt]
/**
 * Redirect 策略：替换为仿真 I/O
 * 适用场景：接收数据函数、轮询状态函数
 *
 * 实现要点：
 * 1. 使用 LCMHal 虚拟设备接口替代寄存器访问
 * 2. 对于固定长度接收，使用 HAL_BE_In
 * 3. 对于可变长度接收，使用 HAL_BE_ENET_ReadFrame
 * 4. 保留超时参数处理逻辑
 */

// 原始代码示例
HAL_StatusTypeDef HAL_UART_Receive(UART_HandleTypeDef *huart,
                                    uint8_t *pData, uint16_t Size,
                                    uint32_t Timeout) {
    uint32_t tickstart = HAL_GetTick();
    for (int i = 0; i < Size; i++) {
        // 轮询接收标志
        while (!(huart->Instance->SR & USART_SR_RXNE)) {
            // 超时检查
            if ((HAL_GetTick() - tickstart) > Timeout) {
                return HAL_TIMEOUT;
            }
        }
        // [REDIRECT] 寄存器读取将被替换
        pData[i] = (uint8_t)(huart->Instance->DR & 0xFF);
    }
    huart->RxXferCount = Size;
    return HAL_OK;
}

// Redirect 策略替换代码
HAL_StatusTypeDef HAL_UART_Receive(UART_HandleTypeDef *huart,
                                    uint8_t *pData, uint16_t Size,
                                    uint32_t Timeout) {
    // [REDIRECT] 使用虚拟设备接口替代寄存器读取
    return HAL_BE_In(pData, Size);  // 从虚拟设备读取数据
    
    // 保留状态更新（如果原函数有）
    huart->RxXferCount = Size;
}
\end{lstlisting}

\subsubsection{Event 策略实现}

Event 策略适用于中断处理函数，引入事件触发机制：

\begin{lstlisting}[style=CStyle,abovecaptionskip=0pt,belowcaptionskip=0pt]
/**
 * Event 策略：事件驱动模拟
 * 适用场景：中断处理函数（IRQHandler）
 *
 * 实现要点：
 * 1. 移除硬件寄存器写入（中断标志清除）
 * 2. 保留中断标志检查逻辑
 * 3. 保留 OS 相关调用（信号量、任务通知）
 * 4. 确保数据处理的执行路径可达
 */

// 原始代码示例
void UART0_IRQHandler(void) {
    if ((UART0->S1 & UART_S1_RDRF_MASK) != 0U) {
        uint8_t data = UART0->D;
        process_received_data(data);
        
        // 通知 RTOS 任务
        BaseType_t xHigherPriorityTaskWoken = pdFALSE;
        xSemaphoreGiveFromISR(semaphore, &xHigherPriorityTaskWoken);
        portYIELD_FROM_ISR(xHigherPriorityTaskWoken);
    }
    // [EVENT] 清除中断标志
    UART0->S1 |= UART_S1_RDRF_MASK;
}

// Event 策略替换代码
void UART0_IRQHandler(void) {
    // [EVENT] 模拟中断标志被设置（从虚拟设备接收数据）
    uint8_t data;
    HAL_BE_In(&data, 1);  // 从虚拟设备接收数据
    
    if ((UART0->S1 & UART_S1_RDRF_MASK) != 0U) {
        process_received_data(data);
        
        // 保留 OS 相关调用
        BaseType_t xHigherPriorityTaskWoken = pdFALSE;
        xSemaphoreGiveFromISR(semaphore, &xHigherPriorityTaskWoken);
        portYIELD_FROM_ISR(xHigherPriorityTaskWoken);
    }
    // [EVENT] 无需清除硬件中断标志，由仿真器管理
}
\end{lstlisting}

\subsection{编译集成实现}

\subsubsection{Builder Agent 核心流程}

Builder Agent 负责将替换代码集成回工程并执行编译。其核心流程如下：

\begin{lstlisting}[style=Python,abovecaptionskip=0pt,belowcaptionskip=0pt]
class BuilderAgent:
    """项目构建智能体"""
    
    def __init__(self, script_path: str, project_path: str):
        self.script_path = script_path
        self.project_path = project_path
        self._graph = None
    
    async def build_graph(self):
        """构建 Builder Agent 工作流"""
        if self._graph is not None:
            return self._graph
        
        # MCP 客户端连接 CodeQL 服务器
        client = MultiServerMCPClient({...})
        
        # 获取工具：编译、文件写入、函数信息查询等
        tools = await client.get_tools()
        
        # 定义工作流节点
        async def build_node(state: AgentState):
            """构建节点：执行编译"""
            messages = state["messages"]
            response = await model_with_structured_output.ainvoke(messages)
            return {"messages": [response], "final_response": response}
        
        # 构建图
        workflow = StateGraph(AgentState)
        workflow.add_node("build", build_node)
        workflow.add_edge(START, "build")
        workflow.add_edge("build", END)
        
        self._graph = workflow.compile()
        return self._graph
    
    async def run(self, build_config: dict) -> BuildOutput:
        """执行构建流程"""
        graph = await self.build_graph()
        initial_state = {
            "messages": [
                {"role": "system", "content": system_prompting_en},
                {"role": "user", "content": f"Build project: {build_config}"}
            ]
        }
        
        result = await graph.ainvoke(initial_state)
        return result["final_response"]
\end{lstlisting}

\subsubsection{错误解析器实现}

错误解析器负责解析编译器输出，提取错误类型、位置和原因：

\begin{lstlisting}[style=Python,abovecaptionskip=0pt,belowcaptionskip=0pt]
class BuildErrorParser:
    """编译错误解析器"""
    
    # GCC 错误格式正则表达式
    GCC_ERROR_PATTERN = re.compile(
        r'(?P<file>.*?):(?P<line>\d+):(?P<col>\d+):\s*'
        r'(?P<type>error|warning):\s*'
        r'(?P<message>.*?)$'
    )
    
    # Clang 错误格式正则表达式
    CLANG_ERROR_PATTERN = re.compile(
        r'(?P<file>.*?):(?P<line>\d+):(?P<col>\d+):\s*'
        r'(?P<type>error|warning):\s*'
        r'(?P<message>.*?)$'
    )
    
    def parse(self, stderr: str) -> List[BuildError]:
        """解析编译错误输出"""
        errors = []
        lines = stderr.split('\n')
        
        for line in lines:
            # 尝试 GCC 格式
            match = self.GCC_ERROR_PATTERN.match(line)
            if not match:
                # 尝试 Clang 格式
                match = self.CLANG_ERROR_PATTERN.match(line)
            
            if match:
                error = BuildError(
                    file=match.group('file'),
                    line=int(match.group('line')),
                    col=int(match.group('col')),
                    type=match.group('type'),
                    message=match.group('message')
                )
                errors.append(error)
        
        return errors
    
    def classify_error(self, error: BuildError) -> str:
        """错误分类"""
        if "undeclared" in error.message:
            return "SYNTAX"
        elif "incompatible" in error.message:
            return "TYPE"
        elif "undefined reference" in error.message:
            return "LINK"
        else:
            return "OTHER"
\end{lstlisting}

\subsubsection{自动修复循环}

自动修复循环是编译集成的关键机制，其伪代码如下：

\begin{lstlisting}[style=Python,abovecaptionskip=0pt,belowcaptionskip=0pt]
Algorithm: Auto-Fix Loop
Input: function_name, replacement_code, max_retries=3

for retry in range(max_retries):
    // 尝试编译
    build_result = compile_project()
    
    if build_result.success:
        // 编译成功，退出循环
        return SUCCESS
    
    // 编译失败，解析错误
    errors = parse_build_errors(build_result.stderr)
    
    // 按错误类型分类
    for error in errors:
        error_type = classify_error(error)
        
        if error_type == "SYNTAX":
            // 语法错误：调用 Fixer Agent 修复
            fix_result = await FixerAgent.fix(
                function_name, error, replacement_code
            )
        elif error_type == "TYPE":
            // 类型错误：调整类型定义
            fix_result = adjust_types(error, replacement_code)
        else:
            // 其他错误：标记为需要人工检查
            return NEEDCHECK
    
    // 应用修复后的代码
    replacement_code = fix_result.fixed_code
    
    // 继续下一轮编译
    if retry == max_retries - 1:
        // 达到重试上限，失败
        return FAILED

return SUCCESS
\end{lstlisting}

\section{本章小结（实现部分）}
\label{sec:implementation_summary_impl}

本章详细介绍了基于驱动程序替换的固件仿真系统的具体实现。系统的实现主要特点包括：

\textbf{（1）模块化架构}：各组件通过标准接口交互，便于扩展和维护

\textbf{（2）LangGraph 智能体}：基于 LangGraph 实现了具有状态管理能力的智能体工作流，支持工具调用和条件分支

\textbf{（3）严格的代码生成规则}：通过 code\_generation\_rules.md 确保生成的代码只使用工具返回或项目头文件中存在的符号

\textbf{（4）完整的验证机制}：包括 Rubric 检查、编译验证和动态测试，形成多层次的代码质量保障

\textbf{（5）自动错误修复}：通过 Fixer Agent 实现编译错误的自动检测和修复，减少人工干预

这些实现细节与第三章的设计紧密对应，验证了系统设计的可行性和有效性。
